\section{直觉理解与改进}

\subsection{按奖励加权的模仿学习}

策略梯度实际上是\textbf{按轨迹奖励加权的模仿学习}：增强高 reward 的动作，抑制低 reward 的动作。相比纯监督，这种写法允许我们通过 stochastic policy 做探索。

\begin{knowledgebox}{试错学习的本质}
\begin{itemize}
    \item 策略生成采样，reward 决定 gradient 的方向；
    \item 更灵活地处理稀疏 reward，因为采样中包含了整个轨迹；
    \item 优化目标是期望 reward，而不是单步准确率。
\end{itemize}
\end{knowledgebox}

\subsection{Reward-to-Go 与因果性}

我们希望某个 time step 的梯度只受该 step 及之后 reward 的影响。用 reward-to-go 可以实现因果性：
$$\nabla_\theta J(\theta) \approx \frac{1}{N}\sum_{i=1}^{N}\sum_{t=1}^{T} \nabla_\theta \log \pi_\theta(a_{i,t} \mid s_{i,t}) \left(\sum_{t'=t}^{T} r(s_{i,t'}, a_{i,t'}) \right).$$
该表达式在每个动作处只观测未来 reward，所以更符合 causal effect。

\begin{table}[H]
\centering
\small
\begin{tabular}{p{0.3\textwidth}p{0.45\textwidth}p{0.2\textwidth}}
\toprule
术语 & 含义 & 作用 \\
\midrule
Trajectory reward & 整个轨迹累加 reward & 全局 supervision \\
Reward-to-go & 当前时刻到终点的 reward & 保持因果性 \\
Baseline & 常数或估值函数 & 减少 variance \\
\bottomrule
\end{tabular}
\caption{策略梯度估计中常用的 reward 表达}
\end{table}

\subsection{Baseline 与 Advantage}

\begin{warningbox}{对奖励尺度敏感}
如果所有奖励都是正的，梯度会增加所有动作的似然——即使有些动作只是“没那么差”。这会显著增加梯度方差。
\end{warningbox}

通过从 reward-to-go 中减去 baseline $b(s_t)$，我们不改变期望但可以降低方差：
$$\nabla_\theta J(\theta) = \mathbb{E}_{\tau \sim p_\theta(\tau)} \left[\sum_t \nabla_\theta \log \pi_\theta(a_t \mid s_t) \cdot (G_t - b(s_t))\right],$$
其中 $G_t$ 是 reward-to-go。若 $b(s_t)=V^\pi(s_t)$，得到 advantage estimator $A^\pi(s_t,a_t)$。

\begin{importantbox}{Advantage 的直觉}
Advantage 衡量某动作好于当前状态平均水平的程度。正 advantage 强化该动作，负 advantage 抑制。
\end{importantbox}

\subsection{代理目标与自动微分}

将 reward-to-go 和 baseline 写入 surrogate objective：
$$\tilde{J}(\theta) = \frac{1}{N}\sum_{i=1}^{N}\sum_{t=1}^{T} \log \pi_\theta(a_{i,t} \mid s_{i,t}) \cdot \text{adv}_{i,t}.$$
该形式便于用自动微分计算梯度，并且允许每条数据都同时参与多步优化，而不必单独反向传播 $N \times T$ 次。

\subsection{本章小结}

本章从直觉角度解析策略梯度：重点在于按奖励加权的模仿、因果性的 reward-to-go，以及 baseline/advantage 的 variance 减少作用。代理目标函数使得梯度计算可以以标准方式在框架中实现。

